\section{Introduction}
\label{sec:intro}

Foundation models pre-trained at scale on medical images or image--text pairs are an increasingly common
starting point for medical image analysis. Self-supervised encoders for
radiology~\cite{perezgarcia2024scalable}, computational pathology~\cite{chen2024generalpurpose,xu2024wholeslide,vorontsov2024model} and
ophthalmology~\cite{zhou2023model} --- most of them trained by the DINOv2
recipe~\cite{oquab2024dinov2}, which is also the general comparator most of the studies reviewed here use
--- are reused frozen across dozens of tasks;
dual-encoder vision--language models (VLMs) built on the CLIP recipe~\cite{radford2021learning}, such as
BiomedCLIP~\cite{zhang2023biomedclip} and CONCH~\cite{lu2024visuallanguage}, support zero-shot
classification and retrieval; multimodal large
language models (MLLMs) such as LLaVA-Med~\cite{li2023llavamed} and MedGemma~\cite{sellergren2025medgemma}
answer image questions; and MAIRA-2~\cite{bannur2024maira2} drafts reports
(Table~\ref{tab:models}, Fig.~\ref{fig:overview}). Reuse
across institutions, modalities and tasks is what makes them foundation models~\cite{bommasani2021opportunities}.

This reuse makes \emph{what internal representations encode} consequential. Self-reported race decodable
from medical images~\cite{banerjee2021ai} and from a frozen foundation-model encoder~\cite{zheng2024demographic}
is available, in principle, to any downstream model that reads those features; a chest-radiograph vision--language system can answer accurately
without the image~\cite{lotfinia2026vision}; and a clinically named concept present in a
representation need not be used~\cite{basu2026interpretability}. Regulatory guidance asks device
developers to declare bias sources and system dependencies~\cite{fda2025aiguidance,euaiact2024}.
Input-pixel attribution, dominant in medical explainable AI, does not characterise feature content and fails
faithfulness criteria in medical-image protocols~\cite{arun2021trustworthiness,sadanandan2026attention},
whereas feature-space interpretability targets the representations themselves.

Linear probing~\cite{alain2016understanding}, concept activation vectors~\cite{kim2018interpretability},
network dissection~\cite{bau2017network} and representational similarity analysis~\cite{kriegeskorte2008representational}
came from convolutional networks and language models and were applied to medical classifiers~\cite{lucieri2020interpretability}.
Since 2023, work has expanded to medical image--text embedding geometry~\cite{restrepo2026cone},
demographic probing of foundation-model embeddings~\cite{zheng2024demographic}, sparse autoencoders
for pathology, radiology and mammography encoders and medical-MLLM residual streams~\cite{bouzid2025insights,nakka2025mammosae,wesp2026sparse},
and steering, editing and patching of medical multimodal models~\cite{nooralahzadeh2026universal,zhu2026evaluating,sadanandan2026mechanistically},
most of it from 2025 onwards (Fig.~\ref{fig:corpus}a) and spread across machine-learning, medical-imaging and clinical
venues.

Existing surveys cover medical XAI, medical foundation models, or general interpretability separately; none
reviews representation-level evidence across medical foundation models (Table~\ref{tab:surveys}).

\textbf{Scope.} This survey covers measurements of and interventions on internal representations in
medical vision, vision--language and multimodal foundation models, including their language
decoders, and the general-domain methods they build on. Clinical language models, foundation models for
electronic health records (EHR), physiological signals and omics, and generative medical image models
fall outside the core scope and enter only as marked anchors or instruments; input attribution enters
only when it tests a representation-level claim (\S\ref{sec:method}). Terminology follows \S\ref{sec:prelim}.

\textbf{Contributions.} The contributions are fourfold. \textit{First}, a taxonomy of feature-space
interpretability methods by methodology, locus and concept provenance, with boundary-case rules
(\S\ref{sec:taxonomy}). It is applied through record-level coding of \Ncore{} core studies, one record per
study, family, evidence type and locus (\S\ref{sec:method}); \Npreprint{} preprints are among them, \Nrepstudies{} are returned by a
logged, replayable retrieval, and the supplement lists every record
(Tables~S1--S6) and every finding's evidence profile (Tables~S9--S11). \textit{Second}, a family-by-family
review of what each family can and cannot establish, and of their requirements, maturity and combinations
(\S\ref{sec:methods}). \textit{Third}, consolidated datasets, concept resources, benchmarks and protocols,
with an assessment of evidence comparability and study quality (\S\ref{sec:evaluation}); a
protocol-bounded evidence map of that corpus, \Fnrows{} cross-family statements in three classes, each
read against \Vnclasses{} validity-test classes
that state what a test licenses beyond the design's base reading and what it does not
(\S\ref{sec:findings}, Table~\ref{tab:vcov}); and a review of experimental uses
(\S\ref{sec:applications}). \textit{Fourth}, open problems and what each solution must demonstrate
(\S\ref{sec:future}).

\begin{figure*}[!tbp]
\centering
\includegraphics[width=0.98\linewidth]{fig1_overview_draft.pdf}
\caption{Overview. Medical vision and multimodal foundation models share a computational path from a
vision encoder, through an optional connector and language model, to a readout. Six loci mark where a
representation can be taken (\S\ref{sec:loci}). Five method families (\S\ref{sec:methods}) interrogate
these representations. Evidence type is a property of each claim rather than of a family: most
claims in the first four families are observational (what is encoded), most claims in the fifth are
interventional (what is used), and a study may contribute both.}
\label{fig:overview}
\end{figure*}
